\documentclass[literature_review.tex]{subfiles}

\begin{document}
\begin{figure}[H]
    \centering
    \includegraphics{Windows_2000_architecture.png}
    \caption{Windows NT Architecture\\ credit: Grm wnr at English Wikipedia.Later versions were uploaded by Xyzzy n at en.wikipedia. - Transferred from en.wikipedia to Commons., CC BY-SA 3.0, https://commons.wikimedia.org/w/index.php?curid=2242280}
    \label{Windows NT Architecture}
    
\end{figure}
Many modern operating systems are designed with a mixture of the above design philosophies. This includes Windows NT which implements certain operating system components in user space including \texttt{environment subsystems}, \texttt{server processes}, \texttt{system support processes} and \texttt{NTDLL.DLL} \TODO{cite something and etc. also do better desc of windows.}. 
and the XNU kernel which is based on the Mach microkernel as the core, including the \gls{ipc} infrastructure of Mach as a key component of the \texttt{kernel environment}, with the BSD layer being part of userland providing the application interface \citep{Apple:Kernel_Prog}. 

\end{document}